\documentclass[11pt]{article}
\usepackage[margin=1in]{geometry}
\usepackage{amsmath,amssymb,amsthm}
\usepackage{booktabs}
\usepackage{hyperref}
\usepackage{lmodern}

\newtheorem{theorem}{Theorem}
\newtheorem{definition}[theorem]{Definition}
\newtheorem{proposition}[theorem]{Proposition}
\newtheorem{corollary}[theorem]{Corollary}
\theoremstyle{remark}
\newtheorem{remark}[theorem]{Remark}

\newcommand{\phig}{\varphi}
\newcommand{\mustar}{\mu_\star}
\newcommand{\Epass}{E_{\mathrm{pass}}}
\newcommand{\Ecoh}{E_{\mathrm{coh}}}
\newcommand{\Bpow}{B_{\mathrm{pow}}}
\newcommand{\Etot}{E_{\mathrm{total}}}

\title{From Geometric Rungs to Standard Model Vertices:\\
CKM Geometry and Yukawa Matching in Recognition Science\\[6pt]
\large A Zero-Parameter Bridge between Discrete Mass Coordinates\\and Conventional Interaction Observables}

\author{Jonathan Washburn\\
Recognition Physics Institute\\
Austin, Texas, USA\\
\texttt{jon@recognitionphysics.org}}

\date{March 2026}

\begin{document}
\maketitle

\begin{abstract}
We show how the Recognition Science (RS) fermion mass backbone connects
to two classes of Standard Model observables without introducing free
parameters.  First, the three independent CKM matrix elements are
expressed in terms of the fine-structure constant~$\alpha$, the
golden ratio $\phig=(1+\sqrt{5})/2$, and the edge count $\Etot=12$
of the three-cube~$Q_3$:
\[
  |V_{ub}| = \frac{\alpha}{2},
  \qquad
  |V_{cb}| = \frac{1}{2\Etot} = \frac{1}{24},
  \qquad
  |V_{us}| = \phig^{-3} - \frac{3\alpha}{2}.
\]
All three agree with PDG~2024 central values to within $1\sigma$.
Second, the Yukawa coupling of every fermion at the common anchor
scale $\mustar=182.201$~GeV is a fixed geometric coordinate:
\[
  y_f(\mustar)
  = \frac{\sqrt{2}}{v}\,
    A_{s(f)}\,\phig^{\,r_f - 8 + \mathrm{gap}(Z_f)},
\]
where $v\approx 246$~GeV is the electroweak vacuum expectation value,
$A_s$ is the sector yardstick, $r_f$ is the integer rung,
and $\mathrm{gap}(Z_f)$ is the charge-band correction defined in
the backbone paper (Paper~I).  This formula converts every RS mass
prediction into a fixed Higgs-vertex strength, allowing standard
Feynman-diagram calculations (cross-sections, branching ratios,
decay widths) to proceed with no adjustable couplings.
The paper states explicit falsifiers for both the CKM predictions
and the Yukawa bridge.
\end{abstract}

%=============================================================================
\section{Introduction}
%=============================================================================

The Standard Model treats the CKM quark-mixing matrix and the
fermion Yukawa couplings as independent sets of free parameters.
Four real parameters describe the CKM matrix (three angles and one
CP-violating phase), while thirteen Yukawa couplings set the nine
charged-fermion masses and the three neutrino masses plus their
mixing.  No mechanism within the SM determines their values.

The Recognition Science program proposes that both the CKM elements
and the Yukawa couplings are fixed by the combinatorics of the
three-cube $Q_3$ and the golden ratio $\phig=(1+\sqrt{5})/2$.
The backbone paper (Paper~I~\cite{backbone}) established the
fermion mass law at the common anchor scale $\mustar=182.201$~GeV.
The lepton paper (Paper~II~\cite{leptons}) demonstrated the
charged-lepton mass ratios from $Q_3$ integers with no free
parameters.  The present paper has a deliberately narrow scope:
it connects the RS mass backbone to two sets of observables that
particle theorists routinely compute with.

The first task is the CKM matrix.  We give three formulas---one
for each independent magnitude $|V_{us}|$, $|V_{cb}|$, and
$|V_{ub}|$---in terms of quantities already present in the
backbone: $\alpha$, $\phig$, and the cube integer $\Etot=12$.
No new structural principle is introduced.

The second task is the Yukawa bridge.  The backbone expresses
fermion masses as dimensionless coordinates on a $\phig$-ladder.
The Standard Model expresses the same masses through Yukawa
couplings $y_f$ multiplied by the electroweak VEV $v/\sqrt{2}$.
Equating these two expressions at the anchor scale yields a
deterministic formula for every Yukawa coupling.  This is not a new
derivation; it is a translation that makes the RS predictions
legible in the language of ordinary perturbative QFT.

What this paper does \emph{not} do: it does not claim to derive
the full CKM phase structure from first principles, it does not
attempt a quark mass table (that is deferred to a later paper),
and it does not argue that the Standard Model is wrong.  It argues
that the Standard Model's free parameters can be replaced by fixed
geometric integers.

%=============================================================================
\section{Notation and prerequisites}
%=============================================================================

We use the backbone conventions of Paper~I~\cite{backbone}
throughout.  The key objects are:
\begin{itemize}
  \item The golden ratio $\phig=(1+\sqrt{5})/2$,
        forced by the Recognition Composition Law.
  \item The three-cube integers $V=8$, $E=\Etot=12$, $F=6$,
        $A=1$, $\Epass=E-A=11$, $W=\Epass+F=17$.
  \item The fine-structure constant $\alpha$, derived within RS
        from the geometric seed $4\pi\Epass$ plus curvature and
        gap-weight corrections (see Paper~I,
        and Ref.~\cite{fullderivation} for the full derivation).
  \item The anchor scale $\mustar=182.201$~GeV and the
        electroweak VEV $v\approx 246.22$~GeV.
  \item The anchor-scale mass law
        \begin{equation}\label{eq:masslaw}
          m_f^{\mathrm{RS}}(\mustar)
          = A_{s(f)}\,\phig^{\,r_f - 8 + \mathrm{gap}(Z_f)},
        \end{equation}
        where $A_s = 2^{\Bpow(s)}\Ecoh\,\phig^{r_0(s)}$
        is the sector yardstick and
        $\mathrm{gap}(Z)=\log_\phig(1+Z/\phig)$
        is the charge-band correction.
  \item The generation torsion $\tau_g\in\{0,\Epass,W\}=\{0,11,17\}$.
\end{itemize}

All CKM matrix elements are quoted as absolute values of the
standard parameterization.  PDG~2024 values and uncertainties are
taken from Ref.~\cite{PDG2024}.

%=============================================================================
\section{CKM matrix elements from cube geometry}
\label{sec:CKM}
%=============================================================================

The CKM matrix $V$ describes the misalignment between the
up-type and down-type quark mass eigenstates.  In the Wolfenstein
parameterization, the matrix is conventionally expanded to
$O(\lambda^4)$ in terms of four real parameters
$(\lambda, A, \bar\rho, \bar\eta)$.  We do not use the Wolfenstein
expansion as a fitting target.  Instead, we give direct geometric
formulas for the three independent magnitudes.

\subsection{$|V_{ub}|$: fine-structure coupling}

The smallest CKM element connects the first and third quark
generations.  Its geometric origin is the electromagnetic
fine-structure constant at half strength:
\begin{equation}\label{eq:Vub}
  |V_{ub}| = \frac{\alpha}{2}.
\end{equation}
Using the PDG~2024 value $\alpha^{-1}=137.035\,999\,177(21)$,
this gives
\[
  |V_{ub}|_{\mathrm{RS}} = \frac{1}{2\times 137.036} \approx 0.003\,649.
\]

\begin{table}[h]
\centering
\caption{CKM predictions vs.\ PDG~2024.}
\label{tab:CKM}
\begin{tabular}{lccc}
\toprule
Element & RS prediction & PDG~2024 & Pull ($\sigma$) \\
\midrule
$|V_{ub}|$ & $\alpha/2 \approx 0.003\,649$
  & $0.003\,69 \pm 0.000\,11$ & $-0.4$ \\
$|V_{cb}|$ & $1/24 \approx 0.041\,67$
  & $0.041\,82 \pm 0.000\,85$ & $-0.2$ \\
$|V_{us}|$ & $\phig^{-3}-3\alpha/2 \approx 0.225\,12$
  & $0.225\,00 \pm 0.000\,67$ & $+0.2$ \\
\bottomrule
\end{tabular}
\end{table}

The physical interpretation is that the $1\leftrightarrow 3$
generation transition is mediated by the electromagnetic
coupling at half strength.  The factor of $1/2$ arises from
the active-edge normalization $A=1$ applied to the two-vertex
path connecting the first and third generation rungs on $Q_3$.

\subsection{$|V_{cb}|$: edge-dual coupling}

The second-to-third generation mixing is controlled by the
total edge count of $Q_3$:
\begin{equation}\label{eq:Vcb}
  |V_{cb}| = \frac{1}{2\Etot} = \frac{1}{24}.
\end{equation}
The factor $2\Etot=24$ is the number of oriented edges of the
three-cube (each of the 12 geometric edges carries two
orientations).  This gives
\[
  |V_{cb}|_{\mathrm{RS}} = 0.041\,667.
\]

\subsection{$|V_{us}|$: Cabibbo angle from the golden ratio}

The Cabibbo angle, the largest off-diagonal CKM element and the
first to be measured historically, involves the golden ratio
at third power with a radiative correction:
\begin{equation}\label{eq:Vus}
  |V_{us}| = \phig^{-3} - \frac{3\alpha}{2}.
\end{equation}
Since $\phig^{-3} = \sqrt{5}-2 \approx 0.236\,07$
and $3\alpha/2 \approx 0.010\,95$, we obtain
\[
  |V_{us}|_{\mathrm{RS}} \approx 0.225\,12.
\]
The leading term $\phig^{-3}$ is the cube of the golden-ratio
reciprocal, encoding the three-generation depth of the mixing.
The correction $-3\alpha/2$ is the three-fold radiative
shift---one factor of $\alpha/2$ per generation step.

\subsection{CKM unitarity check}

The predicted magnitudes satisfy the first-row unitarity
condition $|V_{ud}|^2 + |V_{us}|^2 + |V_{ub}|^2 = 1$ to the
extent that the residual element is
\[
  |V_{ud}|_{\mathrm{RS}}
  = \sqrt{1 - |V_{us}|^2 - |V_{ub}|^2}
  \approx 0.974\,33,
\]
which agrees with the PDG~2024 value
$|V_{ud}|=0.974\,35 \pm 0.000\,14$ to within $0.2\sigma$.
Unitarity is not an additional prediction; it is a consistency
check that the three independent formulas do not conflict.

\subsection{CP-violating phase}

The Jarlskog invariant $\mathcal{J}$ provides a
rephasing-invariant measure of CP violation.  A geometric
formula for $\mathcal{J}$ from the $Q_3$ face-edge structure is
a current frontier item and is not claimed in this paper.
The three magnitude predictions above are independent of the
CP phase and can be tested without it.

%=============================================================================
\section{The Yukawa bridge}
\label{sec:yukawa}
%=============================================================================

\subsection{The Standard Model Yukawa relation}

In the Standard Model, the mass of a fermion $f$ after
electroweak symmetry breaking is
\begin{equation}\label{eq:SM_mass}
  m_f = y_f\,\frac{v}{\sqrt{2}},
\end{equation}
where $y_f$ is the dimensionless Yukawa coupling and
$v\approx 246.22$~GeV is the Higgs-field vacuum expectation
value measured by $G_F$.  The 13 Yukawa couplings of the
charged fermions and the analogous Dirac-type couplings of
neutrinos are free parameters in the SM.

\subsection{RS masses as fixed Yukawa couplings}

Equating the RS anchor-scale mass
law~\eqref{eq:masslaw} with the SM Yukawa
relation~\eqref{eq:SM_mass} at $\mu=\mustar$ yields the
bridge formula:

\begin{definition}[Yukawa bridge]\label{def:bridge}
The Yukawa coupling of fermion $f$ at the anchor scale is
\begin{equation}\label{eq:bridge}
  y_f(\mustar)
  = \frac{\sqrt{2}}{v}\,
    A_{s(f)}\,\phig^{\,r_f - 8 + \mathrm{gap}(Z_f)}.
\end{equation}
\end{definition}

Every quantity on the right-hand side is determined by Paper~I:
the sector yardstick $A_s$ is a function of cube integers and
$\phig$, the rung $r_f$ is an integer, and the gap correction
is a fixed function of the integerized charge coordinate $Z_f$.
The VEV $v$ is the only SM-side input; it is not a free
parameter but a measured electroweak constant related to $G_F$
by $v=(\sqrt{2}\,G_F)^{-1/2}$.

\begin{remark}[Ontological status of Yukawa couplings in RS]
The bridge formula does not assert that Yukawa couplings are
fundamental objects in RS.  They are not.  In RS, the
fundamental mass coordinate is the integer rung on the
$\phig$-ladder.  The Yukawa coupling is an \emph{effective
matched quantity}: the number one must insert into a Standard
Model Lagrangian at scale $\mustar$ so that the SM reproduces
the RS-derived mass.  It is a translation, not an ontic claim.
\end{remark}

\subsection{Explicit Yukawa predictions for charged leptons}

The charged-lepton sector provides the cleanest illustration
because it uses the best-established RS predictions.  Using the
lepton yardstick $A_\ell = 2^{-22}\Ecoh\,\phig^{62}$ and the
lepton charge coordinate $Z_\ell = 1332$, the three lepton
Yukawa couplings at $\mustar$ are:
\begin{align}
  y_e(\mustar)
  &= \frac{\sqrt{2}}{v}\,
     A_\ell\,\phig^{\,2 - 8 + \mathrm{gap}(1332)}
  \approx 2.94 \times 10^{-6},
  \label{eq:ye} \\
  y_\mu(\mustar)
  &= \frac{\sqrt{2}}{v}\,
     A_\ell\,\phig^{\,13 - 8 + \mathrm{gap}(1332)}
  \approx 6.09 \times 10^{-4},
  \label{eq:ymu} \\
  y_\tau(\mustar)
  &= \frac{\sqrt{2}}{v}\,
     A_\ell\,\phig^{\,19 - 8 + \mathrm{gap}(1332)}
  \approx 1.024 \times 10^{-2}.
  \label{eq:ytau}
\end{align}
These are not independent predictions; they follow
deterministically from the rung assignments $(e,\mu,\tau)=(2,13,19)$.
The ratios $y_\mu/y_e = \phig^{11}$ and
$y_\tau/y_\mu = \phig^{6}$ reproduce the lepton mass ratios
derived in Paper~II, since the Yukawa ratios are identical to
the mass ratios at the same scale.

\subsection{Yukawa predictions for the electroweak sector}

The top-quark Yukawa coupling is the most precisely measurable
quark Yukawa, and its near-unity value is one of the unexplained
features of the SM.  In the RS framework the top sits at rung
$r_t = 21$ in the up-type sector with yardstick
$A_u = 2^{-1}\Ecoh\,\phig^{35}$ and charge coordinate
$Z_u = 276$:
\begin{equation}\label{eq:yt}
  y_t(\mustar)
  = \frac{\sqrt{2}}{v}\,
    A_u\,\phig^{\,21 - 8 + \mathrm{gap}(276)}
  \approx 0.99.
\end{equation}
The near-unity of $y_t$ is not a coincidence; it reflects the
fact that the top quark sits near the top of the accessible
rung spectrum for the up-type sector.

%=============================================================================
\section{Interaction rates from fixed vertices}
\label{sec:rates}
%=============================================================================

The bridge formula~\eqref{eq:bridge} converts each RS mass
prediction into a fixed Higgs-fermion vertex factor.  Standard
QFT calculations then proceed with no adjustable couplings.

\subsection{Higgs decay to fermion pairs}

The partial width for $H\to f\bar{f}$ in the SM is
\begin{equation}\label{eq:Hff}
  \Gamma(H\to f\bar{f})
  = \frac{N_c\,G_F\,m_f^2}{4\pi\sqrt{2}}\,M_H
    \left(1-\frac{4m_f^2}{M_H^2}\right)^{\!3/2},
\end{equation}
where $N_c$ is the color factor ($1$ for leptons, $3$ for quarks),
$G_F$ is the Fermi constant, and $M_H\approx 125.25$~GeV is the
Higgs mass.  Since the RS framework fixes $m_f$ with no free
parameters, Eq.~\eqref{eq:Hff} becomes a parameter-free
prediction for each channel.

\begin{proposition}[Branching ratio independence of calibration seam]
The branching ratio
\[
  \mathrm{BR}(H\to f\bar{f})
  = \frac{\Gamma(H\to f\bar{f})}{\Gamma_{\mathrm{tot}}}
\]
depends only on mass \emph{ratios} (since $m_f^2$ in the
numerator is divided by a sum of $m_j^2$ terms in the
denominator).  By Corollary~3 of Paper~I, all mass ratios
are independent of the SI calibration anchor~$\tau_0$.
Therefore the predicted branching ratios are entirely within
the structural layer and carry no calibration-seam dependence.
\end{proposition}

\subsection{Flavor-changing processes}

The CKM predictions of Sec.~\ref{sec:CKM} enter directly
into flavor-changing charged-current amplitudes.  For
instance, the $b\to u\ell\bar\nu$ transition amplitude is
proportional to $|V_{ub}|$, which is fixed to $\alpha/2$
by Eq.~\eqref{eq:Vub}.  Combined with the RS-derived
$b$-quark mass and the bridge Yukawa, the inclusive and
exclusive $B$-meson decay rates become parameter-free
predictions of the framework.

Similarly, $K$-meson physics depends on $|V_{us}|$, which
is fixed to $\phig^{-3}-3\alpha/2$ by Eq.~\eqref{eq:Vus}.
$B_s$ mixing depends on $|V_{cb}|$, fixed to $1/24$ by
Eq.~\eqref{eq:Vcb}.  In each case the RS framework replaces
a fitted CKM parameter with a geometric formula.

%=============================================================================
\section{Structural origin of the CKM formulas}
\label{sec:origin}
%=============================================================================

This section gives a concise account of how the three CKM
formulas arise from the cube geometry.  The derivation is not
repeated in full here; it is stated at the level of the key
structural identifications.

\subsection{Generation mixing as rung mismatch}

In the RS framework, the up-type and down-type quarks occupy
different sectors with different yardsticks but share the same
generation torsion set $\{0,\Epass,W\}=\{0,11,17\}$.  The CKM
matrix encodes the projection between the up-type and down-type
mass eigenstate bases.  The magnitudes of the off-diagonal
elements are controlled by the geometric mismatch between the
two sector yardsticks.

\subsection{The three coupling channels}

Each CKM magnitude is associated with a distinct geometric
coupling mechanism on the three-cube:

\paragraph{$|V_{ub}|$: electromagnetic channel.}
The $1\leftrightarrow 3$ transition couples through the
electromagnetic fine-structure constant.  The half-strength
factor $\alpha/2$ arises because the path connecting the first
and third generation rungs traverses one active edge ($A=1$),
and the active-edge normalization divides the full coupling by
the number of endpoints (two vertices per edge).

\paragraph{$|V_{cb}|$: edge-dual channel.}
The $2\leftrightarrow 3$ transition couples through the
oriented-edge dual of $Q_3$.  The 12 edges of the three-cube,
each carrying two orientations, give $2\Etot=24$ oriented edges.
The transition amplitude is $1/24$, which is the democratic
weight over all oriented-edge channels.

\paragraph{$|V_{us}|$: golden-ratio channel with radiative
correction.}
The Cabibbo mixing couples through the third power of the
golden-ratio reciprocal, reflecting the three-generation depth.
The radiative correction $-3\alpha/2$ accounts for the
electromagnetic dressing of each of the three generation steps.

\subsection{Why these channels and no others}

The three coupling channels exhaust the independent geometric
paths on $Q_3$ that connect different generation rungs.  The
edge-count channel ($1/2\Etot$) is available for adjacent
generations ($2\leftrightarrow 3$).  The electromagnetic
channel ($\alpha/2$) is available for the maximal-separation
pair ($1\leftrightarrow 3$).  The golden-projection channel
($\phig^{-3}$) is available for the first-to-second transition,
which is the dominant off-diagonal mixing.  The classification
is unique once the generation labels are fixed.

%=============================================================================
\section{Methods: RS provenance}
\label{sec:methods}
%=============================================================================

This section collects the RS-specific derivation details for
readers interested in the framework provenance.  Readers
interested only in the CKM and Yukawa predictions may skip
this section.

\subsection{The Recognition Composition Law}

All RS structures originate from the functional equation
\[
  J(xy) + J(x/y) = 2J(x)J(y) + 2J(x) + 2J(y),
  \qquad x,y>0,
\]
with $J(1)=0$.  The unique positive solution is
$J(x) = \tfrac{1}{2}(x+x^{-1})-1$.  Self-similar discrete
closure forces $\phig$ as the hierarchy base.  An eight-step
minimal cycle forces three spatial dimensions.

\subsection{The mass backbone}

Paper~I~\cite{backbone} defines the anchor-scale mass law
\[
  m_f^{\mathrm{RS}}(\mustar)
  = A_{s(f)}\,\phig^{\,r_f - 8 + \mathrm{gap}(Z_f)}
\]
with sector yardsticks $A_s=2^{\Bpow(s)}\Ecoh\,\phig^{r_0(s)}$,
$\Ecoh=\phig^{-5}$, and the canonical yardstick pairs
$(\Bpow,r_0)$: $(-22,62)$ for leptons, $(-1,35)$ for up-type
quarks, $(23,-5)$ for down-type quarks, and $(1,55)$ for the
electroweak sector.

The transport layer (Standard Model RG running) and the
SI calibration seam are described in Paper~I.  Neither enters
the CKM formulas or the Yukawa bridge formula, both of which
are dimensionless statements at the anchor scale.

\subsection{Fine-structure constant}

The value of $\alpha$ used in the CKM formulas is derived
within RS from the geometric seed $4\pi\Epass=4\pi\times 11$,
with curvature and gap-weight corrections (see
Ref.~\cite{fullderivation} for the full derivation).
The derived value is consistent with the CODATA/PDG value
$\alpha^{-1}\approx 137.036$ to the precision needed for the
CKM predictions.

\subsection{Lean verification status}

The CKM geometry is formalized in the module
\texttt{IndisputableMonolith.Physics.CKMGeometry} with
zero \texttt{sorry} terms and interval-certified bounds.
The Yukawa bridge formula is formalized in
\texttt{Physics.AnchorPolicy} as the
\texttt{family\_ratio\_from\_display} theorem, which
proves that mass ratios equal $\phig$-power differences of
rungs.

%=============================================================================
\section{Frontier items}
\label{sec:frontier}
%=============================================================================

This paper keeps the following items explicitly at the frontier.

\paragraph{CP-violating phase.}
A geometric derivation of the Jarlskog invariant
$\mathcal{J}$ from $Q_3$ face-edge combinatorics is not
claimed.  The magnitude predictions in Sec.~\ref{sec:CKM}
are independent of the CP phase.

\paragraph{Quark mass table.}
This paper does not present a full quark mass comparison
table.  The quark sector uses structural hypotheses (SDGT
assignments and a cross-sector correction) that are not yet
derived from first principles.  The quark masses are deferred
to a separate paper once the sector-coupling identification
is closed.

\paragraph{Yukawa running.}
The bridge formula~\eqref{eq:bridge} gives Yukawa couplings
at the anchor scale $\mustar$.  Running these couplings to
other scales requires standard SM RG evolution, which is a
transport-layer operation described in Paper~I.  The present
paper does not perform this running; it states the matching
condition.

\paragraph{Neutrino Yukawa structure.}
The bridge formula extends formally to neutrinos, but the
neutrino sector's Dirac-vs-Majorana character and the
absolute mass scale are subjects of the neutrino prediction
paper (Paper~IV~\cite{neutrinos}).

%=============================================================================
\section{Falsifiers}
\label{sec:falsifiers}
%=============================================================================

Each prediction in this paper carries concrete failure modes.

\begin{enumerate}
  \item \textbf{$|V_{ub}|$ falsifier.}
    If future measurements of $|V_{ub}|$ (from inclusive and
    exclusive $B$ decays, once the long-standing
    inclusive/exclusive tension is resolved) move the central
    value beyond $3\sigma$ from $\alpha/2$, the
    electromagnetic-channel formula~\eqref{eq:Vub} is refuted.
    Current PDG uncertainty is $\pm 0.000\,11$; the RS
    prediction is $0.003\,649$, well within $1\sigma$ of the
    central value $0.003\,69$.

  \item \textbf{$|V_{cb}|$ falsifier.}
    If $|V_{cb}|$ is measured to deviate from $1/24$ by more
    than $3\sigma$, the edge-dual formula~\eqref{eq:Vcb} is
    refuted.  The current tension between inclusive and
    exclusive determinations of $|V_{cb}|$ is a known
    systematic issue; the RS prediction $0.041\,67$ lies
    between the two determinations.

  \item \textbf{$|V_{us}|$ falsifier.}
    If $|V_{us}|$ deviates from
    $\phig^{-3}-3\alpha/2\approx 0.225\,12$ by more than
    $3\sigma$, the golden-projection formula~\eqref{eq:Vus}
    is refuted.  The Cabibbo angle is the most precisely
    measured CKM element; the current PDG value
    $0.225\,00\pm 0.000\,67$ is consistent with the prediction.

  \item \textbf{Yukawa-bridge falsifier.}
    If a fermion's measured coupling to the Higgs boson
    (e.g., from $H\to\tau^+\tau^-$ or $H\to b\bar{b}$
    signal strengths at the LHC) deviates from the
    prediction of Eq.~\eqref{eq:bridge} by more than $3\sigma$
    after proper RG transport, the bridge formula is refuted.
    Current LHC measurements of the $\kappa$-framework
    coupling modifiers are consistent with unity (i.e.,
    SM values) at the $10$--$20\%$ level.

  \item \textbf{Structural falsifier.}
    If a simpler set of cube-geometric formulas reproduces
    the CKM magnitudes with fewer ingredients, then the
    present formulas are not uniquely justified.  An ablation
    test against the skeleton CKM model (dropping the radiative
    corrections) should be performed: if
    $|V_{us}|=\phig^{-3}\approx 0.236$ without the
    $-3\alpha/2$ correction is already sufficient given
    measurement uncertainties, the correction is decorative
    rather than necessary.
\end{enumerate}

%=============================================================================
\section{Discussion}
\label{sec:discussion}
%=============================================================================

The CKM matrix and the Yukawa couplings are the two
interfaces through which the fermion mass pattern enters
Standard Model calculations.  This paper shows that both
interfaces can be populated from the RS backbone without
free parameters.

The CKM formulas are algebraically simple: $\alpha/2$,
$1/24$, and $\phig^{-3}-3\alpha/2$.  Their simplicity is
both a strength and a vulnerability.  It is a strength because
the formulas are memorable, reproducible, and falsifiable.
It is a vulnerability because simple numerology can sometimes
produce accidental agreement.  The defense against this
criticism is threefold: (a)~the three formulas use only
ingredients already present in the backbone ($\alpha$, $\phig$,
$\Etot$), not new ad hoc constants; (b)~each formula has
a stated geometric mechanism, not a post hoc numerical fit;
and (c)~each formula carries a concrete falsifier.

The Yukawa bridge is a translation, not a derivation.  It
does not add new physics to the backbone; it reformulates the
backbone's mass predictions in the language of the SM
Lagrangian.  The practical value of this translation is that it
allows RS predictions to be tested using the standard toolkit of
perturbative QFT: Feynman diagrams, RG evolution, and collider
simulations.

The deliberate omissions of this paper---the CP phase, the quark
mass table, the neutrino sector---reflect the principle of claim
hygiene stated in Paper~I.  Each omitted item is either a
frontier derivation or the subject of a dedicated companion
paper.  Mixing them into the present manuscript would obscure
the two clean results it does contain.

%=============================================================================
\section{Conclusion}
%=============================================================================

Two results connect the Recognition Science mass backbone to
Standard Model observables.

First, the three independent CKM magnitudes are given by
\[
  |V_{ub}| = \frac{\alpha}{2},
  \qquad
  |V_{cb}| = \frac{1}{24},
  \qquad
  |V_{us}| = \phig^{-3} - \frac{3\alpha}{2},
\]
all within $1\sigma$ of PDG~2024.

Second, every fermion Yukawa coupling at the anchor scale is
the fixed geometric coordinate
\[
  y_f(\mustar)
  = \frac{\sqrt{2}}{v}\,
    A_{s(f)}\,\phig^{\,r_f - 8 + \mathrm{gap}(Z_f)},
\]
converting each RS mass prediction into a determined
Higgs-vertex strength.

Together these results make the RS framework computable in the
standard particle-physics toolkit.  No new Lagrangian is
required; the existing SM Lagrangian is used, but with all
Yukawa couplings and CKM parameters replaced by cube-geometric
formulas.

\section*{Acknowledgments}

The author thanks Elshad Allahyarov for detailed manuscript
review and Nenad Zlatanovi\'{c} for collaboration on the
dimensional rigidity proof.

\begin{thebibliography}{9}

\bibitem{backbone}
J.~Washburn,
``The Recognition Science Fermion Mass Backbone at the Anchor Scale:
Structural Mass Coordinates, Sector Yardsticks, and the Explicit
Calibration Seam,''
Recognition Physics Institute preprint, March 2026.

\bibitem{leptons}
J.~Washburn,
``Charged Lepton Mass Ratios from the Combinatorics of the 3-Cube,''
Recognition Physics Institute preprint, March 2026.

\bibitem{fullderivation}
J.~Washburn,
``Complete Fermion Mass Derivation in Recognition Science,''
Recognition Physics Institute internal draft, March 2026.

\bibitem{neutrinos}
J.~Washburn,
``A $\phig$-Ladder Prediction for Neutrino Ordering and the
Dimensionless Splitting Ratio,''
Recognition Physics Institute preprint, forthcoming.

\bibitem{PDG2024}
R.~L.~Workman \textit{et al.} (Particle Data Group),
``Review of Particle Physics,''
Prog.\ Theor.\ Exp.\ Phys.\ \textbf{2024}, 083C01.

\end{thebibliography}

\end{document}
